% ================================================================== Chapitre 8
\chapter{Piles électrochimiques}
\label{ch:piles}

Une pile convertit l'énergie libérée par une réaction rédox spontanée en énergie électrique. L'enthalpie libre fixe à la fois la tension qu'elle peut délivrer et l'énergie maximale qu'elle peut fournir.

\section{Description d'une pile}

\begin{definition}[pile électrochimique et vocabulaire]
\begin{itemize}
  \item Une \textbf{pile} est un générateur électrique dans lequel une réaction d'oxydoréduction spontanée a lieu \emph{sans contact direct} entre les réactifs : l'oxydation et la réduction se déroulent dans deux \textbf{demi-piles} séparées, reliées par une \textbf{jonction électrolytique} (pont salin, paroi poreuse) qui laisse passer les ions sans mélanger les solutions. Les électrons passent par le circuit extérieur.
  \item L'\textbf{anode} est l'électrode où a lieu l'\emph{oxydation} ; la \textbf{cathode} celle où a lieu la \emph{réduction}. Dans une pile, l'anode est le pôle $\ominus$ et la cathode le pôle $\oplus$.
  \item La \textbf{tension à vide} $e$ (anciennement « force électromotrice ») est la tension $V_\oplus - V_\ominus$ mesurée en circuit ouvert ($i = 0$). Elle vaut $e = E_\oplus - E_\ominus = E_{\text{cathode}} - E_{\text{anode}} > 0$.
\end{itemize}
\end{definition}

\begin{figure}[H]
\centering
\begin{tikzpicture}[>=Stealth,scale=0.95]
  % béchers
  \draw[line width=1pt] (0,0) -- (0,3) (0,0) -- (3,0) -- (3,3);
  \draw[line width=1pt] (6,0) -- (6,3) (6,0) -- (9,0) -- (9,3);
  \fill[c1!10] (0.03,0.03) rectangle (2.97,2.3);
  \fill[c1!25] (6.03,0.03) rectangle (8.97,2.3);
  % électrodes
  \fill[black!45] (1.2,0.6) rectangle (1.6,3.7); \node[font=\small,anchor=east] at (1.15,3.4) {Zn};
  \fill[c2!80!black] (7.4,0.6) rectangle (7.8,3.7); \node[font=\small,anchor=west] at (7.85,3.4) {Cu};
  \node[font=\small,align=center] at (0.75,1.0) {\ce{Zn^{2+}}\\\ce{SO4^{2-}}};
  \node[font=\small,align=center] at (8.35,1.0) {\ce{Cu^{2+}}\\\ce{SO4^{2-}}};
  % pont salin
  \draw[line width=6pt,black!20] (2.3,1.4) -- (2.3,2.9) -- (6.7,2.9) -- (6.7,1.4);
  \node[font=\scriptsize] at (4.5,3.2) {pont salin (\ce{K+}, \ce{NO3-})};
  \draw[->,c3,line width=0.9pt] (4.85,2.45) -- node[below,font=\scriptsize] {\ce{K+}} (5.95,2.45);
  \draw[->,c1,line width=0.9pt] (4.15,2.45) -- node[below,font=\scriptsize] {\ce{NO3-}} (3.05,2.45);
  % circuit extérieur
  \draw[line width=0.9pt] (1.4,3.7) -- (1.4,5.2) -- (3.7,5.2);
  \draw[line width=0.9pt] (5.3,5.2) -- (7.6,5.2) -- (7.6,3.7);
  \draw[line width=0.9pt,fill=white] (3.7,4.95) rectangle (5.3,5.45); \node[font=\small] at (4.5,5.2) {$R$};
  \draw[->,navy,line width=1.1pt] (2.0,5.55) -- (3.4,5.55) node[midway,above,font=\small] {$e^-$};
  \draw[->,c3,line width=1.1pt] (7.0,4.85) -- (5.6,4.85) node[midway,below,font=\small] {$i$};
  % pôles
  \node[font=\large] at (0.9,4.5) {$\ominus$}; \node[font=\large] at (8.1,4.5) {$\oplus$};
  % légendes
  \node[font=\small,align=center] at (1.5,-0.8) {\textbf{anode} : oxydation\\\ce{Zn -> Zn^{2+} + 2e^-}};
  \node[font=\small,align=center] at (7.5,-0.8) {\textbf{cathode} : réduction\\\ce{Cu^{2+} + 2e^- -> Cu}};
\end{tikzpicture}
\caption{Pile Daniell en fonctionnement. Les électrons circulent de l'anode vers la cathode dans le circuit extérieur ; dans le pont salin, les cations migrent vers la cathode et les anions vers l'anode pour maintenir l'électroneutralité.}
\label{fig:daniell}
\end{figure}

\noindent Notation conventionnelle : $\ominus\ \ce{Zn(s)} \,|\, \ce{Zn^{2+}},\ce{SO4^{2-}} \,\|\, \ce{Cu^{2+}},\ce{SO4^{2-}} \,|\, \ce{Cu(s)}\ \oplus$. Réaction de fonctionnement : \ce{Zn(s) + Cu^{2+} -> Zn^{2+} + Cu(s)}, avec $\gamma = 2$.

\section{Bilan énergétique : enthalpie libre et travail électrique}

\begin{resultat}{Inégalité entre variation d'enthalpie libre et travail électrique}
Pour un système évoluant de façon monotherme et monobare ($T = T_{\text{ext}}$, $P = P_{\text{ext}}$ aux états initial et final), recevant un travail électrique $W_{\text{élec}}$ en plus du travail des forces de pression :
\[ \boxed{\Delta G \leq W_{\text{élec}}} \qquad\text{(égalité si la transformation est réversible).} \]
Pour une pile, $W_{\text{élec}} < 0$ : l'énergie électrique \emph{fournie} au circuit vérifie
\[ W_{\text{fournie}} = -W_{\text{élec}} \leq -\Delta G, \qquad\text{donc}\qquad W_{\max} = -\Delta G. \]
\end{resultat}

\begin{demo}
Sur une évolution élémentaire, avec $\delta W = -P\,\dd V + \delta W_{\text{élec}}$ :
\begin{itemize}
  \item premier principe : $\dd U = -P\,\dd V + \delta W_{\text{élec}} + \delta Q$, donc $\dd H = \dd U + P\,\dd V = \delta W_{\text{élec}} + \delta Q$ (à $P$ fixée) ;
  \item deuxième principe : $\dd S = \dfrac{\delta Q}{T} + \delta S_c$ avec $\delta S_c \geq 0$ ;
  \item à $T$ fixée : $\dd G = \dd H - T\,\dd S = \delta W_{\text{élec}} - T\,\delta S_c \leq \delta W_{\text{élec}}$.
\end{itemize}
En intégrant : $\Delta G \leq W_{\text{élec}}$, l'égalité correspondant à $\delta S_c = 0$.
\end{demo}

\noindent L'enthalpie libre porte bien son nom : $-\Delta G$ est la part de l'énergie de la réaction qui est « libre » d'être convertie en travail. Le reste ($-\Delta H + \Delta G = -T\Delta S$ pour une évolution réversible) est échangé sous forme de chaleur. En fonctionnement réel ($i \neq 0$), l'effet Joule dans la pile et les irréversibilités aux électrodes créent de l'entropie : on récupère moins que $-\Delta G$.

\section{Tension à vide et enthalpie libre de réaction}

\begin{resultat}{Tension à vide d'une pile}
\[ \boxed{e = -\frac{\Delta_rG}{\gamma F}} \qquad\qquad \boxed{e^\circ = -\frac{\Delta_rG^\circ}{\gamma F} = E^\circ_\oplus - E^\circ_\ominus} \]
\[ e = e^\circ - \frac{RT}{\gamma F}\ln Q_r \;\underset{\SI{298}{K}}{=}\; e^\circ - \frac{0{,}06}{\gamma}\log Q_r = \frac{0{,}06}{\gamma}\log\frac{K^\circ}{Q_r}. \]
\end{resultat}

\begin{demo}
Quand l'avancement augmente de $\dd\xi$, $\gamma\,\dd\xi$ moles d'électrons traversent le circuit extérieur de $\ominus$ vers $\oplus$ : la charge $\dd q = \gamma F\,\dd\xi$ « traverse » la pile de $\ominus$ vers $\oplus$ en interne. Si $u = V_\oplus - V_\ominus$ est la tension aux bornes, la pile (en convention générateur) fournit $u\,\dd q$, donc reçoit $\delta W_{\text{élec}} = -u\,\gamma F\,\dd\xi$.

\noindent À la limite réversible ($i \to 0$), $u \to e$ et $\dd G = \delta W_{\text{élec}}$. Or, à $T$ et $P$ fixées, $\dd G = \Delta_rG\,\dd\xi$. D'où $\Delta_rG\,\dd\xi = -\gamma F e\,\dd\xi$, soit $e = -\Delta_rG/(\gamma F)$.
\end{demo}

En combinant avec la mesure $e = E_\oplus - E_\ominus$, on obtient $\Delta_rG = -\gamma F(E_\oplus - E_\ominus)$ : c'est la relation admise au chapitre~\ref{ch:redox}, maintenant justifiée.

\begin{exemple}[pile Daniell]
$e^\circ = E^\circ(\ce{Cu^{2+}/Cu}) - E^\circ(\ce{Zn^{2+}/Zn}) = 0{,}34 - (-0{,}76) = \SI{1,10}{V}$, donc $\Delta_rG^\circ = -2\times96485\times1{,}10 = \SI{-212}{kJ.mol^{-1}}$ et $\log K^\circ = \dfrac{2\times1{,}10}{0{,}059} \approx 37$.
\end{exemple}

\section{Évolution de la tension pendant la décharge}

Au cours de la décharge, $Q_r = \dfrac{[\ce{Zn^{2+}}]}{[\ce{Cu^{2+}}]}$ augmente, donc $e$ diminue. La pile est « usée » lorsque le système atteint l'équilibre :
\[ \boxed{\text{pile usée : } e = 0 \;\Longleftrightarrow\; \Delta_rG = 0 \;\Longleftrightarrow\; Q_r = K^\circ} \]

Comme $K^\circ \approx 10^{37}$, $Q_r$ reste très inférieur à $K^\circ$ pendant presque toute la décharge : $\log Q_r$ varie de quelques unités seulement, et $e$ reste quasi constante. $\Delta_rG \approx$ constante, donc $G(\xi)$ est quasi \emph{affine} : elle ne s'incurve qu'à l'extrême fin, quand \ce{Cu^{2+}} est presque épuisé.

\begin{figure}[H]
\centering
\begin{tikzpicture}
\begin{axis}[cours,width=0.47\linewidth,height=0.4\linewidth,xlabel={$q$ (C)},ylabel={$e$ (V)},
  xmin=0,xmax=1000,ymin=0.85,ymax=1.15,ytick={0.9,1.0,1.1},
  yticklabel style={/pgf/number format/fixed,/pgf/number format/precision=2},
  title={\small tension à vide},title style={yshift=-4pt}]
  \addplot[c1] table[x=q,y=E]{data/daniell.dat};
  \addplot[c3,dashed,line width=0.8pt] coordinates {(965,0.85) (965,1.15)};
  \node[c3,font=\scriptsize,anchor=east,align=right] at (axis cs:955,0.92) {$q_{\max} = \SI{965}{C}$};
\end{axis}
\end{tikzpicture}\hfill
\begin{tikzpicture}
\begin{axis}[cours,width=0.47\linewidth,height=0.4\linewidth,xlabel={$\xi$ (mmol)},ylabel={$G - G(0)$ (kJ)},
  xmin=0,xmax=5.2,ymin=-1.15,ymax=0.05,x filter/.expression={x/(2*96.485)},
  title={\small enthalpie libre},title style={yshift=-4pt}]
  \addplot[c1] table[x=q,y=G_kJ]{data/daniell.dat};
  \addplot[black!50,dashed,line width=0.8pt] coordinates {(0,0) (5,-1.061)};
  \node[font=\scriptsize,anchor=south west,align=left] at (axis cs:0.3,-0.75) {pente $\approx\Delta_rG^\circ$\\$= -2Fe^\circ$};
\end{axis}
\end{tikzpicture}
\caption{Décharge réversible d'une pile Daniell ($c_0 = \SI{0,10}{mol.L^{-1}}$ et $V = \SI{50}{mL}$ dans chaque compartiment). La tension passe de $\SI{1,10}{V}$ à $\SI{1,09}{V}$ à mi-décharge et vaut encore $\SI{1,03}{V}$ à 99\,\% ; elle ne chute qu'à l'épuisement de \ce{Cu^{2+}}.}
\label{fig:decharge}
\end{figure}

\section{Capacité d'une pile}

\begin{definition}[capacité]
La \textbf{capacité} d'une pile est la charge électrique totale qu'elle peut débiter jusqu'à épuisement du réactif limitant :
\[ \boxed{q_{\max} = \gamma\,F\,\xi_{\max}} \]
Elle s'exprime en coulombs ou, en pratique, en \textbf{ampères-heures} : $\SI{1}{A.h} = \SI{3600}{C}$.
\end{definition}

\begin{resultat}{Énergie électrique maximale}
Si $e$ reste quasi constante pendant la décharge (cas usuel, $Q_r \ll K^\circ$) :
\[ W_{\max} = -\Delta G = -\xi_{\max}\,\Delta_rG = \gamma F\,\xi_{\max}\,e = q_{\max}\,e = i\,\Delta t\,e. \]
\end{resultat}

\begin{exemple}[capacité apportée par $\SI{1}{g}$ de zinc]
$\xi_{\max} = \dfrac{m}{M} = \dfrac{1}{65{,}4} = \SI{15,3}{mmol}$ (zinc limitant), d'où
\[ q_{\max} = 2\times96485\times15{,}3\times10^{-3} = \SI{2,95e3}{C} = \frac{2950}{3600} = \SI{0,82}{A.h}. \]
L'énergie correspondante est $W_{\max} \approx q_{\max}e^\circ = 2950\times1{,}10 = \SI{3,2}{kJ}$. Sous un courant de $\SI{10}{mA}$, la pile fonctionnerait $\Delta t = q_{\max}/i \approx \SI{82}{h}$.
\end{exemple}

\begin{methode}[bilan d'une pile]
\begin{enumerate}
  \item Identifier les couples, écrire les demi-équations : réduction à la cathode ($\oplus$, plus haut potentiel), oxydation à l'anode ($\ominus$).
  \item Écrire la réaction de fonctionnement et repérer $\gamma$.
  \item Tension à vide : $e = E_\oplus - E_\ominus$ (Nernst), ou $e = -\Delta_rG/(\gamma F)$ avec des données thermodynamiques.
  \item Réactif limitant $\Rightarrow$ $\xi_{\max}$ $\Rightarrow$ capacité $q_{\max} = \gamma F\xi_{\max}$, énergie $W\approx q_{\max}e$, durée $\Delta t = q_{\max}/i$.
\end{enumerate}
\end{methode}

\begin{ordres}
\begin{itemize}
  \item Pile alcaline AA : $e \approx \SI{1,5}{V}$, capacité de l'ordre de $\SI{2,5}{A.h}$, soit environ $\SI{3,5}{W.h}$.
  \item Accumulateur lithium-ion : $\approx\SI{3,7}{V}$ par élément ; batterie de smartphone : de l'ordre de $\SI{4}{A.h}$, soit environ $\SI{15}{W.h}$.
  \item En électrolyte aqueux, la décomposition de l'eau limite en pratique la tension vers $\SI{2}{V}$ (accumulateur au plomb : $\SI{2,1}{V}$) ; les accumulateurs lithium-ion dépassent $\SI{3}{V}$ grâce à un électrolyte organique.
\end{itemize}
\end{ordres}
